\title{Byte Latent Transformer: Patches Scale Better Than Tokens}

\begin{abstract}
We present the Byte Latent Transformer ({\textsc BLT}), a novel byte-level LLM design that, for the first time, achieves comparable large-scale performance to tokenization-based LLMs, while offering notable gains in both inference robustness and efficiency. In {\textsc BLT}, input bytes are grouped into dynamically sized patches, which act as the fundamental computational units. The boundaries of these patches are determined according to the entropy of subsequent bytes, thus dedicating greater computational resources and model attention to regions of higher data complexity. We perform the first controlled scaling analysis of byte-level models under fixed compute, scaling to 8B parameters and processing up to 4T bytes during training. The empirical findings confirm that byte-level models can be efficiently scaled without reliance on pre-defined vocabularies. Both the training and inference processes see enhancements in efficiency—especially when data allows the model to dynamically adopt longer patches—as well as observable gains in reasoning capabilities and generalization to rare cases. In summary, at equivalent inference cost, {\textsc BLT} demonstrates superior scaling properties compared to tokenization-based LLMs, benefiting from the concurrent expansion of both patch size and model capacity.
\end{abstract}

\section{Introduction}
\label{section:intro}

We present the Byte Latent Transformer~(\textbf{BLT}), an architecture that eliminates the need for tokenization by directly modeling sequences of raw bytes. BLT achieves, for the first time, parity in performance with traditional tokenization-dependent models at scale, while notably enhancing efficiency and robustness (\S\ref{sec:robustness}).
Conventional large language models (llms) typically undergo end-to-end training, with tokenization standing as the primary heuristic pre-processing intervention that converts byte sequences into a fixed token inventory. This procedure imposes biases on string compression that result in several well-documented issues, including sensitivity to particular domains and modalities~\citep{dagangetting}, increased vulnerability to input noise~(\S\ref{sec:robustness}), insufficient orthographic understanding~\citep{edman2024cute}, and systemic inequities across languages~\citep{liang2023xlm,petrov2024language,limisiewicz2024myte}.

Tokenization has historically played a crucial role because training large language models on raw byte sequences incurs prohibitive computational expenses, mainly due to the resultant extended sequence lengths~\citep{xue2022byt5}. To address this, previous research has adopted strategies such as using more computationally efficient self-attention mechanisms~\citep{el2020characterbert, clark2022canine} or shifting to architectures that bypass attention entirely~\citep{wang2024mambabyte}~(\S\ref{section:related_work}). Nonetheless, these solutions tend to benefit primarily the training of \textit{small models}. When scaling to larger models, it is the substantial computational load from the extensive feed-forward network layers—applied to each individual byte—rather than the attention component, that constitutes the chief bottleneck for Transformer-based systems.

To optimize computational resource usage, we introduce a flexible and adaptive approach for aggregating bytes into \textit{patches}~(\S\ref{section:patching}) and present a novel model architecture that integrates both byte-level and patch-level representations. In contrast to tokenization, {\textsc BLT} does not rely on a predefined patch vocabulary. Instead, any subset of bytes can be transformed into latent patch embeddings through compact, trainable encoder and decoder components. Our results demonstrate that this strategy yields a \textit{greater} efficiency in compute allocation compared to models based on tokenization.

Tokenization-based large language models assign an equal computational budget to each token. This approach sacrifices potential efficiency for improved performance, as the process of inducing tokens relies on compression techniques that may not align with the actual complexity inherent to prediction tasks. Our architecture is founded on the principle that computation should be distributed adaptively according to demand. As an illustrative case, predicting the final letters of words rarely requires the capabilities of a large transformer, since such predictions are typically straightforward and exhibit low entropy, particularly in contrast to the challenge of selecting the initial word of a sentence. This insight informs the structure of {\textsc BLT}'s model~(\S\ref{section:architecture}), which includes three transformer blocks: two smaller byte-level \textit{local models}, and one more substantial global \textit{latent transformer}~(\autoref{fig:arch_high_level}). 
To effectively group bytes into patches—and thus intelligently control computational allocation—{\textsc BLT} segments the input based on the entropy associated with predicting the next byte, resulting in context-dependent clusters of bytes that exhibit approximately consistent information content.

We introduce the first comprehensive study on flop-controlled scaling for byte-level models, exploring configurations up to 8B parameters and 4T training bytes. Our results demonstrate that it is feasible to train large-scale models directly from bytes, circumventing the need for fixed-vocabulary tokenization. Notably, {\textsc BLT} achieves flop-controlled performance on par with Llama 3 while delivering up to 50\% reduction in inference flops~(\S\ref{section:scaling})\footnote{The model's computational requirements are measured in Floating Point OPerations (flops).}. 

Additionally, our findings indicate that modeling directly on raw bytes yields substantial advantages, particularly in capturing the long-tail distribution of the data. {\textsc BLT} models exhibit greater resilience to input noise when compared to traditional tokenizer-based counterparts and show superior capabilities at the character level, as evidenced in tasks involving orthography, phonology, and machine translation for low-resource languages~(\S\ref{sec:robustness}).

Moreover, {\textsc BLT} models permit joint scaling of both model capacity and patch size within an unchanged inference flop ceiling. Utilizing longer patch sizes tends to decrease computational requirements, enabling reallocation of resources to enlarge the global latent transformer since this component operates less frequently. Through inference-flop controlled scaling experiments~(\autoref{fig:fixed_inference_scaling}), we observe that this architecture achieves markedly improved scaling dynamics relative to conventional tokenization-dependent models.

To summarize, the primary contributions of this work are as follows: 1) We present {\textsc BLT}, a byte latent llm model that adaptively manages computational resources to enhance flop efficiency; 2) We demonstrate flop-controlled training performance on par with Llama 3 up to the 8B parameter scale, and further show that it is possible to exchange slight reductions in evaluation performance for up to 50\% improvements in flop efficiency; 3) The {\textsc BLT} architecture introduces a novel scaling avenue for llms, allowing for expansion of model size while keeping inference costs constant; 4) We show that {\textsc BLT} exhibits increased resilience to noisy inputs and demonstrates an understanding of sub-word data characteristics often overlooked by token-based llms. All training and inference code for {\textsc BLT} is made available at~\url{https://github.com/facebookresearch/blt}.

\section{Patching: From Individual Bytes to Groups of Bytes}
\label{section:patching}

Dividing byte sequences into discrete \textit{patches} enables {\textsc BLT} to modulate computational resources dynamically in response to varying contextual demands. As illustrated in Figure~\ref{fig:patching_types}, there exist multiple strategies for organizing bytes into patches. More rigorously, a patching function $f_p$ divides a byte sequence $\pmb{x}=\{x_i,|i=1,\ldots n\}$ of length $n$ into a set of $m < n$ patches $\pmb{p}=\{p_j|j=1,\ldots,m\}$ by assigning each $x_i$ a value in \{0,1\}, with 1 signaling the initiation of a new patch. The main Transformer in both token-based and patch-based architectures incurs computational cost proportional to the number of main processing steps, corresponding in {\textsc BLT} to the count of patches output by the selected patching function. Hence, the mean length of each patch—termed the \textit{patch size}—plays a critical role in dictating the computational expense for data processing, both in training and inference scenarios, for any given patching method~(\S\ref{section:flops}).

We proceed to describe three specific patching strategies: using a constant byte count per patch~(\S\ref{section:static-patch}), patching according to whitespace boundaries~(\S\ref{section:white-patch}), and adaptive patching utilizing entropy estimates from a lightweight byte-level language model~(\S\ref{section:dyn-patch}). Furthermore, we elaborate on incremental patching approaches and clarify distinctions between patching and tokenization~(\S\ref{section:bpe-patch}).

\subsection{Strided Patching Every K Bytes}
\label{section:static-patch}

A commonly used approach for grouping bytes is to segment them into fixed-size patches of $k$ bytes, as implemented in MegaByte~\citep{yu2023megabyte}. This method, with its fixed stride, offers ease of use in both training and inference, and also allows for simple adjustments to the average patch size, making it straightforward to manage the computational cost in FLOPs. Nonetheless, this static patching approach has notable limitations. Primarily, computational resources are not adaptively assigned to regions requiring greater attention: for example, a transformer step $j$ could be wasted when only predicting whitespace in source code, or may be insufficient when dealing with information-rich sequences such as mathematical content. Additionally, this fixed strategy can result in inconsistent and non-context-aware segmentation of comparable byte sequences—for instance, causing the same word to be partitioned in varying ways.

\subsection{Space Patching}
\label{section:white-patch}

\citet{slagle2024spacebyte} introduces a straightforward but impactful refinement to strided patching by forming new patches immediately after any space-like bytes\footnote{
    Space-like bytes encompass all bytes that are neither Latin letters, digits, nor UTF-8 continuation bytes. Each patch is additionally required to include at least one byte that is not space-like. }, which often correspond to natural linguistic boundaries in a variety of languages. In this technique, referred to as Space patching, each word receives an explicit latent transformer step (i.e., increased computational allocation), enabling consistent patching of words throughout sequences and dedicating computational resources to challenging prediction tasks, which are typically found directly following spaces. To illustrate, when generating a response to the prompt ``Who composed the Magic Flute?\uline{\hspace{.6em}}'', predicting the initial byte of the answer poses a substantially greater challenge than predicting subsequent bytes after ``M''—as the first letter largely constrains the set of likely completions, making the prediction of ``Mozart'' after this point considerably less complex. Nevertheless, space patching faces limitations: it is not universally applicable across all linguistic and domain contexts, and more critically, it does not accommodate adjustable patch sizes. In the following section, we propose an alternative patching strategy inspired by the observation that the initial bytes of words are typically the most unpredictable, while simultaneously allowing for adaptive control over patch sizes.

\subsection{Entropy Patching: Using Next-Byte Entropies from a Small Byte LM}
\label{section:dyn-patch}

Instead of depending on heuristics like whitespace for segmentation, we adopt a data-driven strategy to pinpoint instances where next-byte predictions are highly uncertain. Our method, termed \textit{entropy patching}, leverages entropy measurements to establish the locations of patch boundaries.

We train a small byte-level auto-regressive language model on the training data for {\textsc BLT} and compute next byte entropies under the LM distribution $p_{e}$ over the byte vocabulary $\mathcal{V}$:

\begin{equation}
H(x_i) = \sum_{v \in \mathcal{V}} p_{e}(x_i=v|\pmb{x}_{<i}) \log p_{e}(x_i=v|\pmb{x}_{<i})
\end{equation}

We explore two approaches for detecting patch boundaries based on entropies $H(x_i)$. The first approach selects points exceeding a fixed global entropy threshold, as shown in~\autoref{fig:patching}. The second method detects points where the entropy is notably higher compared to the preceding value. This latter strategy can alternatively be understood as locating points that disrupt an approximately monotonically decreasing pattern of entropy within a patch.

\begin{align*}
    \text{Global Constraint}\hspace{1cm}H(x_t)& > \theta_g\\
    \text{Approx. Monotonic Constraint}\hspace{1cm}H(x_t)& - H(x_{t-1}) > \theta_r
\end{align*}

Patch boundaries are determined through a simple preprocessing process that occurs as part of dataloading. This approach contrasts with~\citet{nawrot-etal-2023-efficient}, who employ a classifier trained to recognize entropy-based patch boundaries. In our experiments~(\S\ref{section:experiments}), we evaluate both approaches for separating bytes characterized by low and high entropy.

\subsection{The Byte-Pair Encoding (BPE) Tokenizer and Incremental Patching}
\label{section:incremental-patching}
\label{section:bpe-patch}

Numerous contemporary llms, such as our chosen baseline Llama 3, rely on subword tokenization methods such as bpe~\citep{gage1994new, sennrich-etal-2016-neural}. In this work, we define ``tokens'' as groups of bytes selected from a \textit{finite} vocabulary established before the commencement of training. In contrast, ``patches'' denote adaptively constructed byte sequences that do not rely on a predetermined vocabulary. One important distinction between these two approaches is that, in the case of tokens, the model lacks direct exposure to the byte-level features underlying each token.

An important advancement introduced by {\textsc BLT} relative to tokenization-based methods lies in its reconfiguration of the balance between vocabulary size and computational requirements. Traditional language models necessitate that a larger vocabulary leads to, on average, tokens encompassing more content, which results in fewer model steps but simultaneously demands a higher-dimensional output space in the model’s final projection layer. This dynamic restricts the degree to which tokenization-based techniques can meaningfully alter both token granularity and inference efficiency. For instance, Llama 3 achieves an increase in average token size—from 3.7 to 4.4 bytes—by expanding the corresponding embedding table to four times the size used in Llama 2.

During generation, {\textsc BLT} must determine if the current position in the byte sequence represents a patch boundary; this assessment is crucial, as it dictates whether additional computation will be performed by the Latent Transformer. Importantly, this decision must be made based solely on the information available at the present step, without relying on any bytes from the remainder of the sequence that have not yet been produced. Therefore, any patching method employed cannot depend on upcoming bytes to guide the partitioning of the byte sequence. More precisely, a patching scheme $f_p$ exhibits the property of incremental patching if it holds that:
$$f_p(\pmb{x}_{<i}) = f_p(\pmb{x})_{<i}$$
The bpe approach does not constitute an incremental patching scheme, because the same byte prefix may yield different tokenizations depending on its succeeding context; as a result, it fails to meet the above criterion\footnote{Introducing a unique delimiter token at patch boundaries can modify bpe into an incremental patching scheme; however, this adjustment leads to increased byte-sequence length.}.

\section{{\textsc BLT} Architecture}
\label{section:architecture}
{\textsc BLT} consists of a substantial global autoregressive language model working on patch-level representations, complemented by two compact local models. One local model transforms sequences of bytes into patches, while the other decodes patch representations back into the original byte sequences~(\autoref{fig:arch_high_level}).

\subsection{Latent Global Transformer Model}

\textit{The Latent Global Transformer} functions as an autoregressive transformer model $\mathcal{G}$, comprising $l_{\mathcal{G}}$ layers, and is responsible for transforming a series of latent input patch representations $p_j$ into corresponding output patch representations $o_j$. In this manuscript, subscripts $j$ and $i$ are consistently used to refer to patches and bytes, respectively. The global model incorporates a block-causal attention mask~\citep{dubey2024llama}, ensuring that attention at each step is limited to all patches up to and including the current one within a document. During both pre-training and inference, this model accounts for most of the computational cost (measured in flops). Consequently, by determining strategically when this model is activated, we can adaptively manage and allocate computation to different segments of the input and output streams depending on their complexity.

\subsection{Local Encoder}
\label{section:local}

\textit{The Local Encoder Model}, represented as $\mathcal{E}$, is designed as a compact transformer consisting of $l_{\mathcal{E}} << l_{\mathcal{G}}$ layers. Its central function is to rapidly convert an incoming sequence of bytes $b_i$ into rich, informative patch representations $p_j$. Diverging from the standard transformer design, this model introduces an additional cross-attention module after every transformer layer, which facilitates the aggregation of byte representations into patch-level features~(\autoref{fig:crossattn}).

Initially, each byte $b_i$ in the input is embedded via a matrix of size $\mathbb{R}^{256 \times h_{\mathcal{E}}}$, yielding a representation $x_i$. Optionally, these embeddings can be enriched by incorporating hash-embeddings~(\S\ref{sec:hash-emb}) as supplementary features. 

The architecture employs a stack of alternating transformer and cross-attention layers~(\S\ref{sec:enc-cross-attn}), culminating in patch representations $p_i$, which serve as input to the global transformer $\mathcal{G}$. The attention mechanism in the transformer layers is governed by a \textit{local block causal} mask; specifically, each byte can access a specified window of $w_{\mathcal{E}}$ previous bytes, a window which may span dynamic patch limits but never traverses distinct document boundaries. The next sections provide further discussion on the structure of the embedding mechanism and the cross-attention module.

\subsubsection{Encoder Hash n-gram Embeddings}
\label{sec:ngram-emb}
\label{sec:hash-emb}
An essential factor in generating resilient and informative representations at each position $i$ is the inclusion of context from earlier bytes. {\textsc BLT} addresses this by representing each byte $b_i$ both in isolation and as a constituent of a byte n-gram. Specifically, for every position $i$, we begin by generating byte-grams

\begin{align}
    g_{i,n}=\{b_{i-n + 1},\ldots, b_i\}
\end{align}

for every byte index $i$ and for $n$ ranging from three to eight.\footnote{We disregard byte-grams with length $n$ or greater in cases where $i < n$.}

We proceed by defining hash $n$-gram embeddings, which utilize a hash function to assign each byte $n$-gram to an entry in a dedicated embedding table $E_{n}^{hash}$ of predetermined size, with separate tables maintained for each $n$ in the set $\{3,4,5,6,7,8\}$~\citep{bai2010learning}. The corresponding hash embedding for a given $n$-gram is summed with the byte’s own embedding, after which the combined embedding undergoes normalization prior to being submitted to the local encoder model as input. The resulting augmented embedding can be formulated as

\begin{align}
e_i &= x_i + \sum_{n = 3,...,8}E_{n}^{hash}(\text{Hash}(g_{i,n})) \\
\text{where, Hash}(g_{i,n}) &= \text{RollPolyHash}(g_{i,n}) \% |E_{n}^{hash}|
\end{align}

To standardize $e_i$, we divide it by the total count of $n$-gram sizes incremented by one, utilizing RollPolyHash as described in Appendix~\ref{appendix:rollpolyhash}. Section~\ref{section:ablations} discusses our analysis of how varying $n$-gram hash embeddings, including different $n$ values and embedding table sizes, influence the trends in scaling laws when computational resources are fixed. Beyond using hash $n$-gram embeddings, we conducted experiments with frequency-based $n$-gram embeddings; further information regarding these experiments is presented in Appendix~\ref{appendix:ngrams}.

\subsubsection{Encoder Multi-Headed Cross-Attention}
\label{sec:enc-cross-attn}
Our implementation is largely based on the input cross-attention component found in the Perceiver architecture~\citep{jaegle2021perceiver}. However, a key distinction is that the latent representations in our approach are dynamically tied to the representations of variable-sized patches instead of a predetermined set of latent vectors~(\autoref{fig:crossattn}), and each latent only processes the bytes associated with its own patch. Within this module, each patch $p_j$ is associated with a query vector. This query vector is produced by first aggregating the byte embeddings that form patch $p_j$, and then applying a linear transformation, $\mathcal{E}_{C} \in \mathbb{R}^{h_{\mathcal{E}} \times (h_{\mathcal{E}} \times U_{\mathcal{E}})}$, where $U_{\mathcal{E}}$ represents the count of encoder cross-attention heads. More formally, letting $f_{\text{bytes}}(p_j)$ stand for the byte sequence for patch $p_j$, we determine

\begin{align}
P_{0,j} &= \mathcal{E}_{C}(f_\text{bytes}((p_j)), f~ \text{is a pooling function} \\
P_l &= P_{l-1} + W_o\left(\text{softmax}\left(\frac{QK^T}{\sqrt{d_k}}\right)V\right) \\
\text{where } Q_j &= W_q(P_{l-1,j}), K_i = W_k(h_{l-1,i}), V_i = W_v(h_{l-1,i}) \\
h_l &= \text{Encoder-Transformer-Layer}_l(h_{l-1})
\end{align}

Here, $P \in \mathbb{R}^{n_p \times h_{\mathcal{G}}}$ denotes the $n_p$ patch-level representations intended for processing by the global module. These representations are initially constructed by aggregating the byte embeddings $e_i$ associated with each patch $p_j$. The matrices $W_q$, $W_k$, $W_v$, and $W_o$ serve as the projection weights for queries, keys, values, and outputs, respectively, where the keys and values are generated by applying linear transformations to byte-level representations $h_i$ coming from the preceding network layer (with $h_i = e_i$ at the first layer). A masking approach tailored for patching ensures that each query $Q_j$ can only interact with those keys and values linked to the bytes within the same patch $j$. Since the architecture uses multi-head attention over $Q$, $K$, and $V$, and since patch representations typically possess a greater dimensionality ($h_{\mathcal{G}}$) than $h_{\mathcal{E}}$, $P_l$ is structured as multiple attention heads, each of size $h_{\mathcal{E}}$ in cross-attention. These multiple heads are later concatenated to match the $h_{\mathcal{G}}$ dimensionality. Furthermore, before applying projections, a pre-LayerNorm is introduced to the queries, keys, and values, and this cross-attention mechanism omits positional embeddings entirely. The output of the cross-attention module is wrapped with a residual connection.

\subsection{Local Decoder} 

The local decoder $\mathcal{D}$, much like the local encoder, employs a compact transformer architecture comprising $l_{\mathcal{D}} << l_{\mathcal{G}}$ layers. Its purpose is to transform a sequence of global patch embeddings $o_j$ into corresponding raw bytes $y_i$. This decoder generates the raw byte sequence based on the history of decoded bytes, utilizing hidden states derived from the local encoder’s output for the input byte sequence. Within its architecture, it alternates between $l_{\mathcal{D}}$ cross-attention and transformer layers. Here, the cross-attention mechanism precedes the transformer block, first mapping patch representations into byte representations; subsequently, the transformer layer of the local decoder processes this sequence of bytes.

\subsubsection{Decoder Multi-headed Cross-Attention} 

In the decoder cross-attention, the roles of the queries and key/values are interchanged i.e. the byte-representations are now the queries, and the patch representations are now the key/values. The initial byte-representations for the cross-attention are initialized as the byte embeddings from the last encoder layer i.e. $h_{l_{\mathcal{E}}}$. The subsequent byte-representations for layer $l$, $d_{l,i}$ are computed as:

\begin{align}
D_0 &= h_{l_{\mathcal{E}}} \\
B_l &= D_{l-1} + W_o\left(\text{softmax}\left(\frac{QK^T}{\sqrt{d_k}}\right)V\right), \\
  &\text{where } Q_i = W_q(d_{l-1,i}), K_i = W_k(\mathcal{D}_C(o_j)), V_i = W_v(\mathcal{D}_C(o_j)) \\
  D_l &= \text{Decoder-Transformer-layer}_l(B_l)
\end{align}

Here, $W_k$ and $W_v$ denote the projection matrices for the keys and values, which are applied through linear transformations in combination with the split operation $\mathcal{D}_C$ to the patch representations $o_j$ obtained from the global model. The matrix $W_q$ serves as the query projection, mapping byte-level representations $d_{l-1}$ from the preceding decoder transformer layer (or from $h_{l_{\mathcal{E}}}$ if in the initial layer). $W_o$ functions as the output projection matrix, resulting in $B \in \mathbb{R}^{h_{\mathcal{D}} \times n_b}$, with $n_b$ denoting the count of output bytes. The subsequent decoder representations $D_l$ are determined by passing the output of the cross-attention block $B$ through a decoder transformer layer. Mirroring the approach in the local encoder cross-attention, we implement multi-head attention, pre-LayerNorm, exclude positional encodings, and incorporate a residual connection surrounding the cross-attention mechanism.

\section{Experimental Setup}
\label{section:experiments}
We conduct rigorously controlled experiments aimed at comparing {\textsc BLT} against models based on tokenization, ensuring that {\textsc BLT} is not afforded any advantage by access to potentially extended sequence lengths.

\subsection{Pre-training Datasets} In all experimental setups presented in this paper, models across all examined scales undergo pre-training on two principal datasets: (1) The Llama 2 dataset~\citep{touvron2023llama}, encompassing 2 trillion tokens sourced from an array of publicly accessible datasets, subsequently curated through extensive cleaning and filtering processes to enhance quality; and (2) {\textsc BLT}-1T: A novel corpus comprising 1 trillion tokens sourced from multiple public repositories, which also incorporates a portion of the pre-training material shared by Datacomp-LM~\citep{li2024datacomp}. The Llama 2 dataset serves as the basis for scaling law investigations—following the methodology established by~\cite{dubey2024llama}—to identify optimal architectural configurations for {\textsc BLT}, while the {\textsc BLT}-1T dataset is utilized for a full pre-training procedure intended to benchmark against Llama 3 on a suite of downstream evaluations. It is important to note that neither collection contains data derived from Meta products or related services. In addition, when conducting baseline studies with tokenizer-based models, we adopt the Llama 3 tokenizer, which features a 128K token vocabulary; this tokenizer provided superior baseline results compared to the Llama 2 tokenizer in our empirical assessments.

\subsection{Entropy Model} For our experiments, the entropy model is constructed as a byte-level language model and is trained on the same data distribution as the full {\textsc BLT} system. Unless stated otherwise, we employ a transformer architecture featuring 100M parameters, consisting of 14 layers, a hidden size of 512, and employing a sliding window attention span of 512 bytes. All other hyperparameters remain consistent with those utilized in our local and global transformer configurations. Further exploration of different architectural scales, receptive field sizes, and model variations is detailed in \autoref{section:ablations}. Notably, if the model’s receptive field is sufficiently narrow, it is possible to compress the learned entropy model into an effective lookup table representation.

\subsection{Entropy Threshold and Equalizing Context Length} For entropy-based patching approaches, we determine a patching threshold to reach a target average \textit{patch size} across the pretraining data mixture. In contrast to standard tokenization, {\textsc BLT} permits arbitrary selection of \textit{patch size}, a choice which substantially affects the amount of context available to the model. To prevent models with larger patch sizes from benefiting from increased context, we fix the expected number of bytes per batch, thereby controlling for average context length. As a result, when employing larger patch sizes, we shorten the sequence length accordingly. Specifically, for Llama 2 data, each context consists of 8k bytes, whereas on the {\textsc BLT}-1T corpus, the context window is expanded to an average of 16k bytes, all while keeping the batch size near 16M bytes on average.

Although the mean batch size remains unchanged, dynamic patching techniques produce varying proportions of bytes relative to patches when assembling data batches. To enhance computational efficiency, our approach to {\textsc BLT} training organizes patch batches to circumvent the need for padding within the computationally intensive latent transformer. Consequently, all batches are standardized to contain an equal number of patches. Throughout training, we pad—or in some cases truncate—byte sequences to fixed lengths of 12k bytes for Llama 2 and 24k bytes for {\textsc BLT}-1T datasets. This measure is employed to mitigate the risk of memory overload resulting from sequences containing disproportionately large patches.

\subsection{Entropy Model Context}
\label{section:entropy-context}
Through experimentation, we observe that the application of entropy patching results in increasingly extensive patches, particularly in structured formats such as multiple choice questions (refer to patching demonstrated on an MMLU sample in \autoref{fig:mmlu_patching}), where the content tends to be highly repetitive. The emergence of these larger patches can be attributed to the reduced entropy present in the recurring segments within the entropy model context. Consequently, for the full-scale implementation of {\textsc BLT}-Entropy using a patch size of 4.5, we reinitialize the entropy context by inserting new lines and adopt the approximate monotonicity constraint, as this approach is less vulnerable to the "entropy drift" that arises from fluctuating context lengths. Notably, this modification influences only the computation of entropies, while the strategy for determining the entropy threshold remains unchanged.

\subsection{FLOPs Estimation} 
\label{section:flops}

Our approach primarily adheres to the methodology described in Chinchilla~\citep{hoffmann2022training} for calculating transformer flop counts. This encompasses the feed-forward layers, qkvo projections within the self-attention mechanism, as well as the operations involved in computing attention scores and the output projection. However, we diverge by considering the input embedding layer as an efficient lookup table rather than a dense matrix multiplication, thus assigning zero flops to this operation. Consistent with established conventions, we approximate the backward pass as incurring twice the computational cost in flops relative to the forward pass.

To compute flops \textit{per byte} for {\textsc BLT} models, we add up the flops for the local encoder transformer, the global latent transformer, and the local decoder transformer, together with the cross attention blocks in the encoder and the decoder:

\begin{align}
    \text{FL}_{\text{{\textsc BLT}}} &= \frac{1}{n_p}\, \text{Transf. FL}(h_{\mathcal{G}}, l_{\mathcal{G}}, m=n_{ctx}/n_p, V=0)\\
    &\quad + \text{Transf. FL}(h_{\mathcal{E}}, l_{\mathcal{E}}, m=w_{\mathcal{E}}, V=0) \\
    &\quad + \text{Transf. FL}(h_{\mathcal{D}}, l_{\mathcal{D}}, m=w_{\mathcal{D}}, V=256) \\ 
    &\quad + \frac{k}{n_p} \cdot \text{Cross Attn. FL}(h_{\mathcal{E}}, l_{\mathcal{E}}, m=n_p, r=n_p/k) \\ 
    &\quad + \text{Cross Attn. FL}(h_{\mathcal{D}}, l_{\mathcal{D}}, m=k, r=k/n_p) 
\end{align}

In this context, $n_{ctx}$ denotes the sequence length measured in bytes, and $n_p$ represents the patch size. The parameter $r$ specifies the ratio between queries and key/values, while $k$ indicates the proportion between patch-dimension and byte-dimension; that is, $k$ reflects the count of local model segments that, when concatenated, produce the complete global model representation (illustrated in \autoref{fig:crossattn} with $k=2$). The vocabulary size for the output projection is denoted by $V$, which is utilized exclusively within the local decoder. The attention mechanism’s context length, denoted $m$, varies based on whether the module operates on the byte-level or patch-level sequence, resulting in different attention configurations. Consequently, we adapt the computation of attention FLOPs for each module to accommodate these differences. Explicit formulas for calculating Transformer-FLOPs as well as Cross-Attention FLOPs are provided in Appendix~\ref{section:flops-appendix}.

\subsection{Bits-Per-Byte Estimation} Perplexity only makes sense in the context of a fixed tokenizer as it is a measure of the uncertainty for each token. When comparing byte and token-level models, following previous work~\citep{xue2022byt5, yu2023megabyte, wang2024mambabyte}, we instead report Bits-Per-Byte (BPB), a tokenizer independent version of perplexity. Specifically:

\begin{align}
    \text{BPB}(x)&= \frac{\mathcal{L}_{CE}(\pmb{x})}{\ln(2)\cdot n_{\text{bytes}}}
\end{align}

Here, the aggregate cross-entropy loss, which quantifies the uncertainty associated with the data sequence $\pmb{x}$, is divided by both the number of bytes comprising $\pmb{x}$ and a normalization constant.

\subsection{Transformer Architecture Hyperparameters} The transformer modules within {\textsc BLT}, encompassing both the local and global branches, are predominantly based on the Llama~3~\citep{dubey2024llama} design. Within the feed-forward sections, we employ the SwiGLU activation function~\citep{swiglu}. For encoding positional information, we utilize rotary positional embeddings (RoPE)~\citep{RoPe}, configured with $\theta=500000$ as recommended in~\citep{xiong2024effective}, and apply them solely within the self-attention modules. RMSNorm~\citep{RMSNorm} is adopted for layer normalization. Across all self-attention modules utilizing conventional attention masks—specifically \textit{block causal} or \textit{fixed-window block causal}—we deploy Flash attention~\citep{dao2022flashattention}, fixing the window size at 512 in the case of fixed-width attention masks. When handling cross-attention mechanisms that necessitate dynamic, patch-specific attention masks, we implement Flex Attention\footnote{\url{https://pytorch.org/blog/flexattention}}, which enables fused attention kernel computation and thus offers notable improvements in training efficiency.

\subsection{{\textsc BLT}-Specific Hyperparameters} To evaluate the performance of {\textsc BLT} models, our experimental approach considers two primary avenues: investigating scaling properties and assessing downstream task performance. For these analyses, we use model configurations at several parameter scales: 400M, 1B, 2B, 4B, and 8B. Detailed architectural hyperparameters for each configuration can be found in Appendix~Table~\ref{tab:arch}. 

Initialization of queries in the local encoder's initial cross-attention layer is performed using max-pooling. Each {\textsc BLT} model utilizes $500,000$ hashes with one hash function, and n-gram sizes range from 3 up to 8. A uniform learning rate of $4e-4$ is adopted for all models, justified by a hyperparameter sweep between $1e-3$ and $1e-4$ on 400M and 1B models, where the same value is found to be optimal for both token and {\textsc BLT} setups. For experiments concerning scaling properties on the Llama-2 dataset, we adopt batch sizes as prescribed in~\cite{dubey2024llama}, using their byte-wise equivalent where necessary. Optimization is performed using AdamW~\citep{adamw}, with $\beta_1$ set at 0.9, $\beta_2$ at 0.95, and $\epsilon$ of $10^{-8}$. The learning rate schedule features a linear warm-up over 2000 steps followed by cosine decay down to zero. Weight decay is set at 0.1, and a global gradient clipping threshold of 1.0 is enforced.

\section{Scaling Trends}
\label{section:scaling}

We offer a comprehensive analysis of scaling patterns observed in byte-level models, providing insights that can guide the continued development of {\textsc BLT} models. Our investigation addresses shortcomings in earlier studies of byte-level approaches by: (a) conducting a comparative assessment under compute-efficient training regimes, (b) training and evaluating 8B-parameter models with substantial data volumes (reaching 1T tokens or 4T bytes) on various downstream benchmarks, and (c) analyzing scaling behavior in scenarios where inference cost is regulated. Subsequent sections will explore particular strengths associated with byte-sequence modeling.

\subsection{Parameter Matched Compute Optimal Scaling Trends}
\label{section:parameter-matched-scaling}
We leverage the Llama 2 dataset to train multiple \textit{compute-optimal} bpe and {\textsc BLT} models, selecting four different model capacities between 1B and 8B parameters. The relationship between the total training flops and language modeling accuracy is visualized on a representative subset drawn from the training data mixture. For the bpe models, we follow the parameter-to-data ratio recommended as optimal by Llama~3~\citep{dubey2024llama}. This approach is designed to maximize language model performance under a specified training computation budget~\citep{hoffmann2022training}, and serves as a strong reference point for evaluating other models. Correspondingly, for every bpe model, we also train a {\textsc BLT} model with identical training data and ensure its Latent Transformer precisely mirrors the size and structure of the bpe Transformer counterpart.

As shown in~\autoref{fig:scaling-compute-opt} (right), {\textsc BLT} models consistently equal or surpass the performance of their bpe equivalents, and this pattern persists across increasing model sizes and computational budgets. To our knowledge, {\textsc BLT} represents the first instance of a byte-level Transformer architecture exhibiting comparable scaling behavior to BPE-based models under compute-optimal conditions. This outcome supports our hypothesis that the ideal relationship between parameter count and training compute established for bpe models is also applicable to {\textsc BLT}, or at least deviates minimally.

Advancements in model architecture as well as dynamic patching methods are essential for achieving scalability trends similar to those seen with BPE tokenization. In~\autoref{fig:scaling-compute-opt} (left), we present a comparison between models utilizing space-patching and Llama 3. Here, we approximate the performance of SpaceByte \citep{slagle2024spacebyte} by employing {\textsc BLT} with space-patching, deliberately excluding n-gram embeddings and cross-attention modules. While SpaceByte demonstrates improvements over Megabyte, there is still a considerable gap in performance relative to Llama 3. In~\autoref{fig:scaling-compute-opt} (right), we highlight the gains attributable to architectural enhancements alongside dynamic patching. At large scales, {\textsc BLT} achieves parity with leading tokenizer-based approaches, including Llama 3.

Furthermore, our results indicate that for tokenizer-based models, selecting the Llama-3 tokenizer yields better performance compared to using the Llama-2 tokenizer when both are applied to identical training datasets.

In summary, the {\textsc BLT} architecture demonstrates performance characteristics intermediate to those of Llama 2 and 3, particularly when much larger patch sizes are employed. The bpe tokenizers for Llama 2 and 3 produce tokens averaging 3.7 and 4.4 bytes in length, respectively. {\textsc BLT}, on the other hand, is capable of achieving analogous scaling behaviors while utilizing mean patch sizes of 6 or even 8 bytes. As inference flop requirements scale inversely with the average patch size, adopting an 8-byte patch size can result in approximately a 50\% reduction in inference compute. Furthermore, as both the model and dataset sizes increase, models trained with larger patch sizes tend to exhibit improved performance. Specifically, the {\textsc BLT} variant with an 8-byte patch size initially underperforms compared to Llama 2 with bpe at the 1B parameter scale, yet surpasses bpe by the 7B scale. These observations indicate that such patch sizes could provide even greater advantages as the model and data scales increase, suggesting the potential feasibility and effectiveness of even larger patch sizes as models and training budgets grow.

\subsection{Beyond Compute Optimal Task Evaluations}
\label{section:evals}
To further investigate scaling behavior, we continue training an 8B {\textsc BLT} model beyond the compute-optimal threshold using the {\textsc BLT}-1T dataset, which is both larger and of higher quality. We then evaluate model performance across a variety of established classification and generation tasks. The chosen evaluations focus on common sense reasoning, broad world knowledge, and code generation competencies:
\paragraph{Classification tasks} consist of ARC-Easy (0-shot)~\citep{Eval_ARC}, Arc-Challenge (0-shot)~\citep{Eval_ARC}, HellaSwag (0-shot)~\citep{Eval_hellaswag}, PIQA (0-shot)~\citep{Eval_piqa}, and MMLU (5-shot)~\citep{hendrycks2020measuring}. For these tasks, we utilize a prompt-based scoring technique by computing the likelihood associated with each possible choice and calculate mean accuracy scores across datasets.
\paragraph{Coding related generation tasks:} We present pass@1 results on MBPP (3-shot)~\citep{Eval_MBPP} and HumanEval (0-shot)~\citep{Eval_HumanEval} to assess the proficiency of large language models in Python code generation.

In \autoref{tab:evals}, we present a comparative analysis of three models, each trained on the {\textsc BLT}-1T dataset: a model leveraging the bpe Llama 3 tokenizer,\footnote{We opt for the Llama 3 tokenizer, featuring a 128k vocabulary, as it outperforms the Llama 2 tokenizer with its 32k vocabulary.} alongside two distinct {\textsc BLT} model variants. The first employs a space-patching approach ({\textsc BLT}-Space), while the second adopts an entropy-driven patching strategy ({\textsc BLT}-Entropy). The entropy model enforces an approximate monotonicity constraint, and its context is reset with new lines, as elaborated in~\autoref{section:entropy-context}. All models were allotted an equivalent computational (flop) budget during training. For {\textsc BLT}-Entropy, we further refine inference behavior by lowering the entropy threshold from 0.6 to 0.1, a modification that enhances performance on downstream tasks but necessitates a greater number of inference steps.

The {\textsc BLT}-Entropy model surpasses the Llama 3 model in performance on four of the seven evaluated tasks, despite both models being trained on equal byte quantities. This advantage is likely attributed to two primary factors: (1) enhanced efficiency in the utilization of training compute enabled by dynamic patching, and (2) the explicit modeling of information at the byte level rather than at the token level.

Conversely, {\textsc BLT}-Space exhibits inferior performance relative to the Llama 3 tokenizer on all tasks except one. However, it manages to substantially lower inference flops owing to its greater average patch size of 6 bytes. For context, models utilizing bpe and entropy-patching approaches maintain an average patch size of roughly 4.5 bytes when evaluated on the training dataset mix. Under identical training resource constraints, the model with a larger patch size is able to process 30\% more data compared to the other two approaches, which may result in {\textsc BLT} deviating further from the compute-optimal regime.

\subsection{Patches Scale Better Than Tokens} {\textsc BLT} models make it possible to increase both the model size and patch size simultaneously without exceeding the original training and inference computational costs or requiring additional training data. The ability to expand patch size without restriction distinguishes patch-based models from traditional fixed-vocabulary token-based models, and enables them to overcome the efficiency limitations inherent to tokens, as outlined in Section~\ref{section:bpe-patch}. Larger patch sizes reduce computational demands, which in turn allows computational resources to be redirected toward scaling up the global latent transformer, since its frequency of operation decreases.

We perform a fixed inference scaling analysis to investigate whether increasing model size—paired with fewer processing steps over larger patches—yields superior performance compared to smaller models executing more steps. Using models with 400 million and 3.6 billion parameters initialized with the Llama 2 tokenizer, we identify flop-equivalent counterparts using the Llama 3 tokenizer as well as {\textsc BLT}-Entropy models, which average patch sizes of 6 and 8 bytes across the training datamix (refer to~\autoref{tab:fixed-inf-params} for further details about the models). For those models utilizing an 8-byte patch size, we employ 3 encoder layers instead of just 1. Each model is then trained under a variety of training flop budgets.

As illustrated in \autoref{fig:fixed_inference_scaling}, {\textsc BLT} architectures exhibit more favorable scaling behaviors in comparison to tokenization-based models when evaluated at equivalent inference flop levels. For both types of models, BPE-based approaches yield superior performance at lower training compute, but {\textsc BLT} overtakes them soon after the compute-optimal point is exceeded. From a practical standpoint, allocating increased resources to the one-time pretraining phase can result in a higher-performing model under a fixed inference compute constraint. An illustrative case is provided by 8B model families, such as Llama 3.1, which underwent pretraining on a dataset two orders of magnitude greater than what would be considered optimal given their model size and computational cost.

The point at which {\textsc BLT} begins to outperform token-based architectures moves somewhat nearer to the compute-optimal threshold as the models increase in the flop class, shifting from approximately 3 times to 2.5 times the optimal compute allocation. Additionally, in this higher flop class, the model configured with patch size 8 demonstrates a more pronounced scaling curve, surpassing other configurations at an earlier stage. As outlined in Section~\ref{section:parameter-matched-scaling}, increasing patch sizes yields performance more closely aligned with BPE models as the scale of the models increases. We partially attribute this to the relative reduction in flops consumed by the byte-level Encoder and Decoder components, whose resource requirements seem to scale more gradually than those of the Latent Transformer. For example, scaling up the overall parameter count from 400M to 8B (a 20-fold increase) only roughly doubles the share of local model parameters within {\textsc BLT}. Notably, because larger patch sizes impact only the flops expended by the patch Latent Transformer—leaving byte-level modules mostly unaffected—this contributes to the observed patterns. Indeed, this phenomenon explains why the size ratio of {\textsc BLT}-Entropy with patch size 8 grew modestly from 1.6x to 1.7x that of the Llama 2 baseline when advancing to a larger model regime.

In conclusion, our investigation into patch-length scaling reveals that the {\textsc BLT} architecture benefits from concurrent increases in both patch size and model capacity, resulting in superior scaling performance. Moreover, these improvements are maintained—and may become more pronounced—as the models grow larger.

\section{Byte Modeling Improves Robustness} Additionally, we assess how resilient {\textsc BLT} is relative to token-based models that do not incorporate byte-level details. We also introduce a method for converting pretrained token-based models to operate at the byte level.
\label{sec:robustness}

\subsection{Character-Level Tasks}

An initial impetus for developing byte-level models was their inherent resistance to input noise at the byte granularity, along with their capacity to recognize the internal structure of tokens—an area where traditional tokenizer-based models exhibit limitations. To empirically assess these characteristics, we conduct supplementary evaluations using benchmarks specifically designed to test both resilience to input perturbations and sensitivity to character-level information, across English, various other languages, as well as digits and phonemes. The corresponding findings are summarized in Table~\ref{tab:char_tasks}.

\paragraph{Noisy Data} To assess and contrast the resilience of tokenizer-based architectures with that of {\textsc BLT}, we generate noised variants of the benchmark classification datasets detailed in Section~\ref{section:evals}. Five character-level noise methods are utilized to systematically alter the input text: (a) \textit{AntSpeak}: The full text is converted to uppercase, with spaces inserted between every character. (b) \textit{Drop}: Random deletion of 10\% of the text's characters. (c) \textit{RandomCase}: A random half of the characters are made uppercase, with the remaining half rendered in lowercase. (d) \textit{Repeat}: Twenty percent of the characters are repeated, each up to a maximum of four consecutive times. (e) \textit{UpperCase}: Every character is transformed into its uppercase equivalent. During evaluation, each noise transformation is applied independently to the prompt, the completion, or both, and the corresponding results are averaged over these conditions. Table \ref{tab:char_tasks} summarizes the performance on the noised HellaSwag~\citep{Eval_hellaswag} dataset, revealing that {\textsc BLT} consistently surpasses tokenizer-based baselines in robustness, averaging an 8-point improvement over comparably-trained models, and even outperforms the Llama 3.1 variant trained on a substantially larger corpus.

\paragraph{Phonology - Grapheme-to-Phoneme (G2P)} We evaluate {\textsc BLT}'s performance in converting a sequence of graphemes—letters or characters in a word—into its corresponding phonemic representation, which reflects the word's pronunciation. Table \ref{tab:char_tasks} displays outcomes from the G2P task conducted with five examples per class, employing Phonology Bench~\citep{suvarna2024phonologybench}. Our results show that {\textsc BLT} surpasses the baseline model based on the Llama 3 1T tokenizer in this scenario.

\paragraph{CUTE}
To evaluate the capacity for character-level understanding, we test {\textsc BLT} on the CUTE benchmark~\citep{edman2024cute}, which encompasses multiple tasks grouped into three primary types: comprehension of composition, recognition of orthographic similarity, and sequence manipulation capabilities. This suite presents considerable difficulty for typical tokenizer-based architectures, as these models often recognize token spellings but have difficulty leveraging this awareness for direct text manipulation. As presented in Table~\ref{tab:char_tasks}, {\textsc BLT}-Entropy surpasses both BPE Llama 3 variants by over 25 points in performance on this benchmark. Notably, our model exhibits near-perfect performance in character manipulation, attaining 99.9\% accuracy across both spelling-related challenges. Such substantial gains, especially given that {\textsc BLT} was trained on just $\frac{1}{16}$ the data used for Llama 3.1, suggest that character-level features are particularly challenging for BPE architectures to internalize. Examples depicted in Figure~\ref{fig:cute_examples} demonstrate instances where the Llama 3 tokenizer model encounters difficulties that our {\textsc BLT} model is able to handle effectively. The only tasks where BPE holds an advantage are word deletion and insertion, which may prove more intricate for a byte-level system to master; however, the performance gap is relatively modest, and it may be more feasible to scale from character-level to word-level manipulation than the reverse. For all experiments, we employ the original evaluation framework and prompts provided via Huggingface. It is possible that BPE-based models could improve with the application of further prompt engineering.

\paragraph{Low Resource Machine Translation} We assess {\textsc BLT} on both directions of translation across twenty-one low-resource languages, encompassing six widely-studied language families and multiple scripts, utilizing the FLORES-101 benchmark~\citep{goyal2022flores}. SentencePiece BLEU scores are presented in Table~\ref{tab:flores}. Our findings reveal that {\textsc BLT} delivers stronger performance than a system leveraging the Llama 3 tokenizer, surpassing it by 2 BLEU points in translation into English, and by 0.5 points for translation out of English. On well-resourced language pairs, {\textsc BLT} matches or exceeds Llama 3’s results. Notably, {\textsc BLT} achieves superior outcomes over Llama 3 in many low-resource language pairs, highlighting the value of byte-based modeling approaches in handling and generalizing to infrequent or unique byte sequences.

\subsection{Training {\textsc BLT} Using Llama 3 Initialization}
\label{sec:distillation}
We investigate a strategy in which {\textsc BLT} models benefit from existing pre-trained, tokenizer-dependent models to achieve more efficient and accelerated convergence during training. Specifically, this is accomplished by transferring the global transformer parameters from a pre-trained Llama 3.1 to initialize {\textsc BLT}. After initialization, the global transformer parameters are finetuned with a learning rate set to one-tenth of that utilized for the local encoder and decoder components, following the recommended number of optimization steps for Llama 3. A performance comparison with a baseline {\textsc BLT} is summarized in Table~\ref{tab:distill}. The results demonstrate that the Llama 3.1-initialized {\textsc BLT} decisively outperforms both the Llama 3 baseline and the original {\textsc BLT} when controlled for computational budget. Furthermore, even relative to our {\textsc BLT}-Entropy configuration reported in Table~\ref{tab:evals}—which was exposed to a considerably larger corpus (1T tokens)—the {\textsc BLT} model initialized from Llama 3.1 obtains higher MMLU scores. This indicates that such an initialization protocol can substantially decrease training compute requirements while preserving or enhancing final task performance.

This approach can be interpreted as reconfiguring tokenizer-dependent architectures into tokenizer-independent ones, thereby transforming a pre-trained LLaMA 3.1 model into a {\textsc BLT} model. For a thorough evaluation, we present both the original LLaMA 3.1—trained on 15T tokens—in Table \ref{tab:distill} and assess its performance relative to that of the {\textsc BLT} obtained from LLaMA 3. Although our model only exhibits a minor reduction in results for MMLU and HumanEval, the decrease is more pronounced on additional benchmarks. These findings indicate that additional research is essential to make full use of the capabilities of the pre-trained model, particularly by refining the composition of training data and tuning relevant hyperparameters.

\section{Ablations and Discussion}
\label{section:ablations}
Here, we present a series of ablation studies that support our decisions regarding the architecture of {\textsc BLT}, as well as the selection of the patching strategy and hyper-parameter configurations used for the 8B-parameter {\textsc BLT} model trained on the {\textsc BLT}-1T dataset.

\paragraph{Entropy Model Hyper-parameters} To investigate how changes in entropy model size and the length of the context window influence scaling behavior, we train byte-level entropy transformer models spanning a range from 1 million to 100 million parameters, utilizing context window lengths that vary between 64 and 512. We generate scaling law curves depicting bpb versus training flops, derived from our $400m$ and $1b$ {\textsc BLT} models that were trained on the Llama-2 dataset, and display these in \autoref{fig:entropy_ablation}. Our results show that scaling the entropy model along both the size and context length axes yields improved performance; however, gains diminish once the model exceeds 50 million parameters.

\paragraph{Types of Patching}

We conduct an ablation study on the four patching methods described in Section~\ref{section:patching}: 1) Strided Patching employing strides of 4 and 6, 2) Patching that aligns with whitespace boundaries, 3) Patching informed by BPE Tokenizer segments as determined by the Llama 3 tokenizer, and 4) Entropy-driven patching utilizing a compact byte-level language model.

Although dynamic patching shortens the actual sequence length, we ensure that the context length remains consistent across all patching methods by controlling the sequence length. During both training and inference, each model is exposed to the same expected number of bytes per sequence, eliminating potential confounding variables related to context size. The outcomes of these ablation experiments are illustrated in Figure~\ref{fig:scaling-compute-opt}. Notably, every alternative patching strategy surpasses the static patching baseline, with space patching achieving performance nearly equivalent to dynamic entropy based patching.

\autoref{tab:evals_ablation} reports the performance of {\textsc BLT} models across several benchmarks, offering a comparison among models utilizing tokenizers, space patching, and entropy-based patching. All models were trained on the Llama 2 dataset, selecting training steps for optimal performance~\citep{dubey2024llama}. Despite the relative simplicity of space patching—which foregoes the use of an online entropy model during training—we observe that the improvements from entropy-based patching, initially identified through scaling analysis (Section~\ref{section:scaling}), persist in these downstream benchmarks. \footnote{Results for space patching were obtained from earlier experiments lacking cross-attention, but equivalent trends were found when cross-attention was included.}

\paragraph{Cross-Attention} 
\autoref{tab:crossattn_ablations_1b} presents an ablation study investigating the impact of introducing cross-attention at different locations within the encoder and decoder components of {\textsc BLT}. For the encoder cross-attention, we explore several strategies for initializing the queries: (1) assigning an identical learned embedding to each global state, (2) utilizing a hash embedding derived from the bytes within the patch, and (3) applying pooling operations to the encoder's hidden state representations of the patch bytes at the specific encoder layer under consideration.

Our results indicate that applying cross-attention within the \textit{decoder} yields the strongest performance. Employing cross-attention in the encoder leads to modest gains, observed primarily when the queries are initialized via pooling. Furthermore, cross-attention proves especially beneficial for the Common-Crawl dataset and when working with larger patch sizes.

\paragraph{n-gram Hash Embeddings} 

We experiment with n-gram hash embedding vocabularies set to 0, 100K, 200K, and 400K, with the outcomes summarized in Table~\ref{tab:ngram_ablations_1b}. The addition of hash embeddings yields improvements across every domain, with the largest benefits observed in Wikipedia and Github—showing a 0.04 bpb gain as opposed to just 0.01 bpb after 15k steps with the 8B model. At the 8B parameter scale, a reduction in hash vocabulary size from 500K to 300K leads to only a marginal 0.001 bpb difference at 15k steps. These results suggest that employing hash embeddings is essential for narrowing the gap between the performance of {\textsc BLT} and that of tokenizer-based systems, yet the incremental gains taper off beyond 300K hashes. Moreover, the results indicate that hash embeddings and cross-attention have largely additive effects, as they provide distinct performance improvements depending on the dataset.

\paragraph{Local Model Hyperparamaters} Table \ref{tab:local_layers} presents an ablation study exploring different configurations for the number of layers within both the local encoder and decoder components. The results indicate that, when used alongside hash n-gram embeddings, {\textsc BLT} achieves strong performance with a highly simplified encoder consisting of only a single layer, while the decoder benefits from being more complex.

\section{Related Work}
\label{section:related_work}

\textbf{Character-Level RNNs:} Since the inception of neural language models~\citep{sutskever2011generating,mikolov2012subword,graves2013generating}, character-level language modeling has drawn significant interest. This is largely due to its inherent ability to handle novel words natively, eliminating the need for conventional back-off strategies. In their work, \cite{kim2016character} introduce an architecture that ingests only character sequences on the input side. By combining convolutional and highway networks as feature extractors and feeding the result into LSTM-based RNNs, their approach achieved results comparable to contemporaneous RNN-based state-of-the-art language models on English, while surpassing them in morphologically complex languages—highlighting another key benefit of character-level LLMs. Building on this, \cite{kenter2018byte} implemented byte-level LSTM architectures for machine comprehension tasks, showing superior results to word-level systems in languages with rich morphology such as Turkish and Russian. Similarly, \cite{zhang2015character} demonstrated that character-based convolutional networks excel on certain text classification benchmarks, outperforming their word-level counterparts. \citet{chung2019hierarchical} advanced this line of research with hierarchical LSTM frameworks equipped with boundary detectors, enabling the automatic discovery of latent textual hierarchies and enhancing character-level modeling effectiveness. Distinct from attention mechanisms, ByteNet by \cite{kalchbrenner2016neural} applies convolutional layers directly to character inputs for neural machine translation.

\textbf{Character-Level Transformers:} The advent of transformers leveraging attention mechanisms~\citep{vaswani2017attention} in conjunction with subword tokenization schemes~\citep{sennrich-etal-2016-neural} brought notable advancements in neural language modeling and performance on various benchmarks. Nevertheless, the use of word and subword units introduces a preset inductive bias regarding the granularity at which the models interpret language. In pursuit of merging the advantages found in transformer-based models with the encouraging progress seen in character-level language modeling, \citet{al2019character} employed substantially deep transformer architectures. By incorporating auxiliary training objectives, they succeeded in developing transformer models surpassing preceding LSTM-based character language models. Despite these advancements, a substantial performance discrepancy persisted in comparison to word-level large language models. Additionally, GPT-2~\citep{radfordgpt2} noted that, especially with extensive datasets such as the 1 Billion Word Benchmark, byte-level language models continued to lag behind word-level counterparts in effectiveness.

Previous work by \cite{Choe2019BridgingTG} established that transformer-based LLMs operating at the byte level can surpass their subword counterparts when matched for parameter count. However, these byte-level models demand substantially greater computational resources and extend training times. In a similar vein, \cite{el2020characterbert} introduced CharFormer, a BERT-based model which constructs word-level representations via convolutions over character embeddings. This model demonstrated domain-specific improvements—particularly in medical applications—but also incurred higher computational costs. The CANINE architecture proposed by \cite{clark2022canine}, with 150M parameters, processes sequences at the character level utilizing a deep stack of transformers analogous to our global model. In addition, CANINE employs both local transformers and strided convolutions to downsample character sequences, and it surpasses the token-based encoder-only baseline mBERT on various multilingual tasks. In another notable advancement, ByT5~\citep{xue2022byt5} pursued byte-level encoder-decoder models that omit patching operations altogether. Though ByT5 displayed enhanced resilience to input noise and comparable performance to tokenized models even with only a quarter of the data, the absence of patching necessitated expensive attention over every byte, resulting in significant computational load. Directly utilizing bytes as modeling units, rather than subwords, increases input sequence lengths and thus challenges scaling such models efficiently. More recently, leveraging the Mamba Architecture~\citep{gu2023mamba}, which allows for a fixed-size memory state across very long contexts, \cite{wang2024mambabyte} trained a byte-level Mamba model—also without patching—and demonstrated its superiority over byte-level transformers of similar scale (350M parameters) in terms of bits-per-byte across several datasets under flop-matched comparisons.

\textbf{Patching-based approaches:} Leveraging patching methods has demonstrated a capacity to mitigate the typically high computational cost—measured in floating point operations—for byte-level LLMs, while still maintaining model performance. Several studies have provided early evidence of the effectiveness of these methods, primarily at relatively small scales in terms of both model size and training data. In their work, \cite{nawrot-etal-2022-hierarchical} utilize static patching for both downsampling and upsampling within their proposed hourglass transformer, which surpasses other byte-level baselines when evaluated at the 150M parameter level. Building upon this, \cite{nawrot-etal-2023-efficient} introduce dynamic patching methods, such as a learned end-to-end boundary-predictor, a boundary-predictor trained with specific tokenizers, and an entropy-driven patching technique akin to {\textsc BLT}. Their findings reveal that these advancements can lead to better performance compared to standard transformers on language modeling tasks at a 40M parameter scale, using approximately 400M tokens. In addition, \cite{lester2024training} examine the effects of training on sequences that have been compressed using arithmetic coding—achieving compression rates beyond those possible with BPE. Employing an equal-info windows strategy, their approach exceeds byte-level baselines on language modeling metrics, although it still lags behind subword-based models.

Our research takes significant motivation from MegaByte~\citep{yu2023megabyte}, a decoder-only causal LLM architecture utilizing fixed static patching to aggregate and concatenate representations for mapping bytes into patches. The approach employs a specialized local model within the decoder to recover bytes from patches. The original work on MegaByte reveals its capability to perform on par with tokenizer-based models at the 1B parameter scale across a corpus of approximately 400B bytes. Throughout all our experiments, we benchmark MegaByte, observing that its reliance on static patching trails behind contemporary state-of-the-art compute-efficient, tokenizer-based methods when evaluated under equal computational constraints. We show that {\textsc BLT} effectively closes this gap. In a parallel observation, \cite{slagle2024spacebyte} recognize similar limitations with MegaByte, and propose augmentations such as patching based on whitespace and analogous space-like bytes in addition to incorporating a local encoder. Their findings indicate superior performance relative to token-based transformers, particularly in domains like Github and arXiv, again at the 1B parameter scale under compute budget restrictions. In our work, we also test this modified approach and demonstrate that additional architectural innovations are necessary to enable byte-level models to scale further and achieve true parity with the latest token-based architectures, exemplified by models like Llama 3.

\section{Limitations and Future Work}

In this study, we adopt the optimal number of training steps identified for Llama~3~\citep{dubey2024llama} to guide our architectural decisions. Nonetheless, since these scaling laws were derived specifically for transformers operating at the BPE level, applying them to {\textsc BLT} may yield less than ideal ratios between data and parameter sizes. The development of scaling laws tailored to {\textsc BLT} is left as future work, and such efforts may reveal even more advantageous scaling behavior for our architecture. Furthermore, a significant portion of our experimental evaluation has been restricted to models with up to 1B parameters. It is conceivable that the best architectural configurations could shift with further scaling to models with 8B parameters or larger, potentially resulting in enhanced performance at greater scales.

Current transformer libraries and codebases are optimized primarily for architectures that utilize tokenizers. Although we provide theoretical experiments with matched floating point operations and employ efficient solutions like FlexAttention for non-standard layers, our implementations might still lag behind tokenizer-based models regarding actual wall-clock performance and could be further improved through additional optimization efforts.

Whereas {\textsc BLT} relies on an independently trained entropy model for patching purposes, an alternative approach worth exploring in future research would be to jointly optimize the patching model in an end-to-end manner. Section \ref{sec:distillation} describes preliminary experiments that demonstrate promising signs for ``byte-ifying'' existing tokenizer-based models—such as Llama 3, trained on upwards of 10T tokens—by directly initializing and keeping the global transformer’s weights fixed. Expanding on this line of investigation could yield approaches that leverage the advantages of bytefying, and potentially achieve performance surpassing that of current tokenizer-based models, all without needing to retrain these large models from the ground up.

\section{Conclusion}
\label{section:conclusion}

This work introduces the Byte Latent Transformer (\textbf{BLT}), an innovative model architecture that challenges the standard reliance on predetermined vocabulary tokenization in large language models. Through a novel, adaptive approach in which the model learns to group bytes into variable-size patches, {\textsc BLT} optimally directs computational effort according to the underlying data’s complexity. This strategy yields notable gains in both computational efficiency and model resilience. Our large-scale experiments indicate that, at scales up to 8B parameters and 4T bytes, {\textsc BLT} attains performance on par with established tokenization-based architectures such as Llama 3, and is capable of achieving up to a 50\% reduction in inference flops in exchange for only marginal drops in standard metrics. In addition, {\textsc BLT} facilitates a new axis for scaling: enabling concurrent expansion of both model and patch sizes, all while maintaining a constant inference cost—a valuable trait for computational settings often encountered in real-world applications. By directly processing byte-level input, {\textsc BLT} enhances the handling of rare or diverse data types, demonstrates improved robustness to noisy inputs, and exhibits a more nuanced comprehension of sub-word patterns. Collectively, our findings highlight {\textsc BLT} as a compelling contender to traditional tokenization schemes, offering a versatile and efficient framework conducive to next-generation language modeling.

\end{document}